\documentclass[a5paper, twoside, symmetric, nohyper]{tufte-book}
\usepackage[portuges]{babel}
\begin{document}
\chapter{Muitas psicologias, um só compromisso}

O termo ``Psicologia'' tem sua origem nas palavras gregas \emph{psyché}
(alma, espírito) e \emph{logos} (estudo, razão, compreensão). A
Biblioteca Virtual em Saúde define ``Psicologia'' como uma ``ciência da
alma, ou da psique, ou da mente, ou do comportamento''. De forma ampla,
a Psicologia visa estudar e compreender os comportamentos, sentimentos,
emoções, vivências e outras questões inerentes à vida de uma pessoa. Ela
lida com o ser humano em toda a sua complexidade, o que demanda que
as/os profissionais também tenham uma visão crítica acerca da sociedade
e dos desafios territoriais. Isso inclui compreender e intervir em
questões crônicas relacionadas à desigualdade e à opressão social e seus
impactos sobre a saúde integral.

A presença da Psicologia em diversos espaços da sociedade,
consequentemente, tem sido cada vez mais requisitada.

É importante ressaltar que não há somente uma abordagem teórica em
Psicologia: são incontáveis as teorias e formas de analisar e explicar
as ações e sentimentos humanos, cada qual baseando-se num princípio e
numa concepção de mundo. As linhas teóricas mais conhecidas atualmente
são a psicanálise, a psicologia comportamental, a terapia
cognitivo-comportamental (TCC), as abordagens humanista e sistêmica, a
gestalt-terapia e a psicologia junguiana --- mas a lista é infindável e
continua crescendo conforme os anos passam. Cada profissional deve
escolher a abordagem teórica que mais se afina com sua visão de mundo e,
a partir dela, embasar sua prática profissional.

No Brasil, a profissão de psicóloga/o foi regulamentada em 1962, por
meio da Lei nº~4.119/1962. Em nosso país, a Psicologia é composta, em
sua maioria, por profissionais mulheres. No estado de São Paulo, por
exemplo, temos, no momento, 128.794 psicólogas e 21.857 psicólogos com
inscrição ativa. Mas onde estão essas/es psicólogas/os?

Ao contrário do que a maioria da população pensa, a Psicologia não se
limita à clássica atuação em consultório clínico. São inúmeras as áreas
em que uma/um psicóloga/o pode atuar: clínica, social, saúde,
organizacional e do trabalho, hospitalar, escolar, jurídica, do esporte,
do tráfego, psicopedagogia, neuropsicologia, avaliação psicológica,
entre tantas outras. Cada profissional tem muito a contribuir com sua
observação atenta e grande capacidade analítica em qualquer que seja o
campo escolhido para atuação.

Ainda neste contexto da atuação de profissionais, é relevante
compreender a distinção entre Psicologia e psicoterapia/terapia. As
atividades de psicoterapia/terapia, embora sejam costumeiramente
desenvolvidas por psicólogas/os, \textbf{não} são privativas da
Psicologia ou das/os psicólogas/os, de modo que outras/os profissionais
não estão impedidas/os de atuar como terapeutas/psicoterapeutas, desde
que não avancem sobre aspectos privativos da Psicologia (como o uso de
testes psicológicos). Há, entretanto, movimentos e projetos de lei ---
acompanhados pelo CFP, por meio de sua assessoria --- que visam
torná-las privativas.

A restrição do exercício de atividades privativas da Psicologia somente
a profissionais com inscrição ativa no respectivo Conselho Regional de
Psicologia é fundamental para melhor garantir a qualificação dos
serviços prestados e a proteção da saúde da população, considerando o
respeito devido por essas/esses profissionais às funções precípuas dos
conselhos de classe: a regulamentação, orientação, fiscalização e
disciplinamento do exercício da profissão.

A Psicologia, como dissemos, é ampla e diversa. Existem princípios e
referências, no entanto, que valem para todas/os as/os profissionais,
como as diretrizes relacionadas aos Direitos Humanos e a atuação baseada
em parâmetros éticos que permitam oferecer um serviço de qualidade
técnica e científica garantida. Essas/es profissionais devem respaldar a
sua atuação em diretrizes legais construídas pelo Sistema Conselhos de
Psicologia, a exemplo do Código de Ética Profissional da/o Psicóloga/o.

A/o psicóloga/o deve estar fundamentada/o técnica e eticamente em sua
atuação. Para uma atuação integral, também é importante que a/o
profissional esteja bem articulada/o com a rede psicossocial, que
oferece um cuidado a pessoas em sofrimento mental e está ligada às
políticas relacionadas ao Sistema Único de Saúde (SUS) e ao Sistema
Único de Assistência Social (Suas).

Outra base da atuação profissional é o respeito ao sigilo e à
privacidade. Sigilo significa manter sob proteção as informações e os
fatos conhecidos por meio da relação profissional, em que estão
implicadas a confiança e a exposição da intimidade da/o usuária/o.
Toda/o psicóloga/o, em seu exercício profissional, está obrigada/o ao
sigilo, sendo este um dos pontos fundamentais do trabalho profissional,
cabendo à/ao psicóloga/o criar condições adequadas para que não seja
violado. De acordo com o Código de Ética, a/o psicóloga/o tem por dever
profissional manter o sigilo e a privacidade das pessoas atendidas, mas
não pode ser conivente com maus-tratos e violações de direitos humanos,
assumindo o compromisso de denunciar essas situações às instâncias
competentes.

Uma/um profissional de Psicologia possui à sua disposição uma
diversidade de técnicas e intervenções, baseadas em estudos e métodos
científicos. Um ponto importante de ser lembrado é que as/os
psicólogas/os, em sua atuação no campo da saúde, muitas vezes trabalham
de forma interdisciplinar --- ou seja, em diálogo com profissionais como
médicas/os, assistentes sociais, enfermeiras/os, entre outras/os,
respeitadas as especificidades de cada atuação (por exemplo: uma/um
psicóloga/o não possui competência legal para prescrever medicamentos;
no entanto, contribui para um entendimento mais integral da saúde, que
considere a história de cada sujeito e os contextos sociais em que está
inserido).

São diversas as possibilidades de atuação de uma/um psicóloga/o, assim
como são diversas as demandas com que essas/es profissionais lidam. Caso
esteja enfrentando desafios relacionados a sofrimento mental, ou mesmo
que esteja em busca de uma escuta especializada, realizada com ética e
respeito aos Direitos Humanos, saiba que a Psicologia está atuante em
diversos setores da sociedade.

\end{document}
